\section{System F confluence and explicit type checking}
\label{sec:p07-system-f-confluence-and-checking}

Confluence and decidable checking concern different inputs. Confluence
applies to the raw annotated terms, including terms that have no typing
derivation or admit infinite reductions. Checking decides a specified
Church-style judgment by its finite syntax. Neither proof below assumes
normalization. Together with the separate normalization theorem, confluence
also gives uniqueness of the normal form of every well-typed term.

We use exactly Sheet~\ref{sheet:system-f}, with the combined compatible
relation \(\to_{\mathsf F}\) and capture-avoiding substitutions of
Definitions~\ref{def:p07-system-f-sn-raw} and
\ref{def:p07-system-f-sn-substitutions}. In particular there are no eta
contractions and no reductions inside types. The parallel-development
method follows the beta fragment of
\cite[Sections~4.3--4.4]{Selinger2013LambdaNotesV2}; the explicit
type-abstraction and type-application cases are supplied here.
The System~F syntax is also presented in
\cite[Sections~8.1--8.4]{Selinger2013LambdaNotesV2} and
\cite[Sections~11.1--11.2]{GirardLafontTaylor2003Web}.

\subsection{Parallel reduction and substitution}
\label{sec:p07-system-f-parallel}

\begin{hbdefinition}{def:p07-system-f-parallel}{An auxiliary parallel relation}
On raw terms modulo alpha-renaming, let \(M\Rrightarrow_{\mathsf F}N\)
be the least relation with the following closure clauses. They define an
auxiliary relation used in the proof; the book's computational steps remain
the two registered beta contractions.
\begin{enumerate}
\item Every variable satisfies \(x\Rrightarrow_{\mathsf F}x\).
\item If \(M\Rrightarrow_{\mathsf F}M'\) and
      \(N\Rrightarrow_{\mathsf F}N'\), then
      \(MN\Rrightarrow_{\mathsf F}M'N'\).
\item If \(P\Rrightarrow_{\mathsf F}P'\), then
      \(\lambda x{:}\alpha.P\Rrightarrow_{\mathsf F}
        \lambda x{:}\alpha.P'\).
\item If \(P\Rrightarrow_{\mathsf F}P'\), then
      \(\Lambda a.P\Rrightarrow_{\mathsf F}\Lambda a.P'\).
\item If \(P\Rrightarrow_{\mathsf F}P'\), then
      \(P[\alpha]\Rrightarrow_{\mathsf F}P'[\alpha]\).
\item If \(P\Rrightarrow_{\mathsf F}P'\) and
      \(N\Rrightarrow_{\mathsf F}N'\), then
      \((\lambda x{:}\alpha.P)N\Rrightarrow_{\mathsf F}P'[N'/x]\).
\item If \(P\Rrightarrow_{\mathsf F}P'\), then
      \((\Lambda a.P)[\alpha]\Rrightarrow_{\mathsf F}P'[\alpha/a]\).
\end{enumerate}
Binders are renamed fresh for the substitution images in clauses~6 and~7.
The relation can contract several nested redexes, and its reflexive cases
allow some redexes to be retained. No transitivity clause is included.
\index{System F!parallel reduction}\index{confluence!System F}
\end{hbdefinition}

\begin{hblemma}{thm:p07-system-f-parallel-closure}
 {proved=annex:p7.system-f.parallel-closure}{The same finite reductions}
For every raw term \(M\), \(M\Rrightarrow_{\mathsf F}M\). Moreover,
\[
 \to_{\mathsf F}\ \subseteq\ \Rrightarrow_{\mathsf F}
 \subseteq\ \to_{\mathsf F}^*.
\]
Consequently their reflexive transitive closures coincide.
\end{hblemma}

\begin{hblemma}{thm:p07-system-f-parallel-substitution}
 {proved=annex:p7.system-f.parallel-substitution}{Parallel substitution, including mixed binders}
For arbitrary raw terms and capture-avoiding type substitution \(\tau\),
\[
 M\Rrightarrow_{\mathsf F}M'
 \quad\Longrightarrow\quad M\tau\Rrightarrow_{\mathsf F}M'\tau.
\]
If additionally \(N\Rrightarrow_{\mathsf F}N'\), then
\[
 M[N/x]\Rrightarrow_{\mathsf F}M'[N'/x].
\]
These assertions permit open terms, open type arguments, erased
occurrences and arbitrarily many duplicated occurrences of \(x\).
\end{hblemma}

\subsection{Complete development and the join}
\label{sec:p07-system-f-development}

\begin{hbdefinition}{def:p07-system-f-development}{Complete development of the original redexes}
Define \(M^\star\) recursively on the original term syntax:
\[
\begin{aligned}
 x^\star&=x,\\
 (\lambda x{:}\alpha.P)^\star&=\lambda x{:}\alpha.P^\star,\\
 (\Lambda a.P)^\star&=\Lambda a.P^\star.
\end{aligned}
\]
For a term application, distinguish its original function constructor:
\[
 (MN)^\star=
 \begin{cases}
 P^\star[N^\star/x],&M=\lambda x{:}\alpha.P,\\
 M^\star N^\star,&M\text{ is not a lambda abstraction}.
 \end{cases}
\]
For a type application use
\[
 (M[\alpha])^\star=
 \begin{cases}
 P^\star[\alpha/a],&M=\Lambda a.P,\\
 M^\star[\alpha],&M\text{ is not a type abstraction}.
 \end{cases}
\]
Types are unchanged except for the displayed capture-avoiding substitutions.
This is a finite structural recursion. A new redex created by developing
a function is retained when the original function was not an abstraction;
\(M^\star\) need not be normal.
\end{hbdefinition}

\begin{hbtheorem}{thm:p07-system-f-development-join}
 {proved=annex:p7.system-f.development-join}{Every parallel reduct reaches the same development}
For raw terms,
\[
 M\Rrightarrow_{\mathsf F}N
       \quad\Longrightarrow\quad N\Rrightarrow_{\mathsf F}M^\star.
\]
In particular the parallel relation has the diamond property: any two
of its one-step reducts have a common one-step parallel reduct, with
reflexive parallel steps allowed.
\end{hbtheorem}

The finite rectangular joining argument in
Annex~\ref{annex:p7.system-f.confluence} proves the existing
Theorem~\ref{thm:p07-system-f-confluence} from this diamond and the
closure lemma. No termination hypothesis is used. Thus raw divergent
terms are included in the confluence statement.

\begin{hbexample}{def:p07-system-f-mixed-overlaps}{The two mixed nesting patterns}
Consider \(((\Lambda a.(\lambda x{:}\alpha.P)Q)[\beta])\).
Contracting the outer type redex first, then the exposed term redex, gives
\(P[\beta/a][Q[\beta/a]/x]\). Contracting the inner term redex
first, then the type redex, gives \((P[Q/x])[\beta/a]\).
These are equal by the type/term substitution identity with capture
avoided.

Conversely consider \((\lambda x{:}\alpha.(\Lambda a.P)[\beta])N\).
Rename \(a\) fresh for the free type variables of \(N\).
The two orders give \((P[N/x])[\beta/a]\) and
\((P[\beta/a])[N/x]\), which agree under that freshness condition.
The parallel-substitution lemma handles reductions duplicated through
\(N\) as well as these displayed root contractions. The development
proof also covers two term redexes, two type redexes, disjoint positions,
and contexts that expose a new redex after reduction.
\end{hbexample}

\subsection{Finite representations and alpha equality of types}
\label{sec:p07-system-f-checking-types}

\begin{hbdefinition}{def:p07-system-f-checking-inputs}{Effective syntax and context validation}
Represent the two disjoint sorts of variable names by finite strings with
decidable equality. A type context \(\Delta\) lists distinct type
variables; a term context \(\Gamma\) assigns one finite type to each
of its distinct term variables. Malformed syntax is rejected.
Bound names may be renamed consistently and capture-free; free names in
\(\Delta\) and \(\Gamma\) are kept fixed.

Write \(\mathsf{WF}(\Delta,\alpha)\) for the formation check:
accept a variable precisely when it occurs in \(\Delta\); accept
an arrow when both components pass; at \(\forall a.\beta\), rename
\(a\) fresh for \(\Delta\) and check \(\beta\) under \(\Delta,a\).
Validate \(\Gamma\) by applying this check to every declared type.

To compare two types modulo bound renaming, traverse each with a stack
of enclosing universal binders. Encode a bound-variable occurrence by
its distance to the nearest matching binder in the stack, and a free
variable by its actual name. Encode an arrow by the ordered pair of its
children and a universal type by a binder marker followed by its encoded
body. Two types are alpha-equal precisely when these finite encodings agree.
\end{hbdefinition}

\begin{hblemma}{thm:p07-system-f-type-syntax-decision}
 {proved=annex:p7.system-f.type-syntax-decision}{Deciding formation, alpha equality and substitution}
The formation check terminates and accepts exactly the types derivable
over \(\Delta\) by the sheet's formation rules. The stack comparison
terminates and decides equality of types modulo bound renaming while
preserving free names. Capture-avoiding type substitution is effective
on this finite syntax and returns a finite type. It preserves formation
under the hypotheses of Lemma~\ref{thm:p07-system-f-sn-typing-substitution}.
\end{hblemma}

\subsection{Synthesizing and checking the explicitly supplied type}
\label{sec:p07-system-f-checking-algorithm}

\begin{hbmethod}{def:p07-system-f-synthesis}{The syntax-directed synthesis procedure}
For validated contexts define \(\mathsf{Synth}(\Delta,\Gamma,M)\)
to return a type or the failure marker \(\mathsf{fail}\).
Before descending through a binder, consistently rename its bound name
fresh for the current contexts and the finitely many free names involved.
This is permitted alpha-renaming of the input, not renaming its free
context declarations.
\begin{enumerate}
\item At \(x\), return its declared type in \(\Gamma\); fail if
      there is no declaration.
\item At \(\lambda x{:}\alpha.P\), require
      \(\mathsf{WF}(\Delta,\alpha)\). Synthesize \(P\) in
      \(\Gamma,x:\alpha\). If its type is \(\beta\), return
      \(\alpha\mathbin{\to_{\rm ty}}\beta\); propagate failure.
\item At \(PN\), synthesize both subterms. Require the first result
      to have arrow shape \(\alpha\mathbin{\to_{\rm ty}}\beta\),
      and the second to be alpha-equal to \(\alpha\). Return \(\beta\)
      if these tests succeed; otherwise fail.
\item At \(\Lambda a.P\), after the fresh renaming require \(a\)
      absent from \(\Delta\) and from the free variables of every type
      in \(\Gamma\). Synthesize \(P\) under \(\Delta,a;\Gamma\).
      Return \(\forall a.\beta\) if that call returns \(\beta\),
      and propagate failure otherwise.
\item At \(P[\alpha]\), require \(\mathsf{WF}(\Delta,\alpha)\)
      and synthesize \(P\). Its result must have universal shape
      \(\forall a.\beta\). Return the capture-avoiding instance
      \(\beta[\alpha/a]\); otherwise fail.
\end{enumerate}
The checker for an input judgment first validates the proposed result
type, then accepts precisely when synthesis succeeds and its result is
alpha-equal to that proposed type. There is no implicit type instantiation,
generalization or search for missing lambda annotations.
\end{hbmethod}

\begin{hbtheorem}{thm:p07-system-f-synthesis-correct}
 {proved=annex:p7.system-f.synthesis-correct}{Termination, exactness and uniqueness of the synthesized type}
For finite validated contexts and a finite explicit term, synthesis
terminates. It returns \(\beta\) exactly when the term has a derivation
\(\Delta;\Gamma\derives_{\mathsf F}M:\beta\), modulo the permitted
bound renamings. If the term is typable, every type assignable in those
same contexts is alpha-equal to the returned type. If synthesis fails,
no such typing derivation exists.
\end{hbtheorem}

This proves the stated Church-style checking theorem
\ref{thm:p07-system-f-church-checking}; its final decision argument is
Annex~\ref{annex:p7.system-f.church-checking}. The recursion follows
proper subterms of the supplied term. Although substitution can enlarge
an output type, it does not add recursive term-checking calls on that
output. This termination argument uses neither confluence nor term
normalization, and asserts no polynomial bound for a chosen representation.

\begin{hbexample}{def:p07-system-f-checking-examples}{Explicit instantiation and fixed free names}
In contexts \(\Delta=b\) and
\(\Gamma=f:\forall a.(a\mathbin{\to_{\rm ty}}a),\ z:b\),
synthesis of \(f[b]z\) returns \(b\): the type application returns
\(b\mathbin{\to_{\rm ty}}b\), and the term application matches
its domain. Synthesis of \(fz\) fails because the type of \(f\) has
universal, not arrow, outer shape. The missing instance is not guessed.

Alpha equality must distinguish free and bound names. Under a context
containing the free type variable \(a\), the types \(\forall b.a\)
and \(\forall b.b\) are different. Their bodies encode a free name
and a bound index, respectively. By contrast \(\forall a.a\) and
\(\forall b.b\) are alpha-equal as types. Consistent freshening of
an input binder must never capture or rename a free context variable.
\end{hbexample}
